%% ============================================================
%%  ChurnIQ – Intelligent Player Churn Prediction
%%  Project Report  |  LaTeX  |  Milestone 1 (Mid-Semester)
%% ============================================================
\documentclass[12pt, a4paper]{article}

%% ─── Packages ───────────────────────────────────────────────
\usepackage[a4paper, top=2.5cm, bottom=2.5cm, left=2.8cm, right=2.8cm]{geometry}
\usepackage{times}
\usepackage{setspace}
\onehalfspacing
\usepackage{graphicx}
\usepackage{float}
\usepackage{booktabs}
\usepackage{array}
\usepackage{multirow}
\usepackage{amsmath}
\usepackage{amssymb}
\usepackage{hyperref}
\usepackage{xcolor}
\usepackage{listings}
\usepackage{caption}
\usepackage{subcaption}
\usepackage{enumitem}
\usepackage{fancyhdr}
\usepackage{titlesec}
\usepackage{parskip}
\usepackage{cite}
\usepackage{tikz}
\usepackage{pgfplots}
\usepackage{tcolorbox}
\usepackage{tabularx}
\pgfplotsset{compat=1.18}
\tcbuselibrary{skins, breakable}

%% ─── Hyperlink colours ──────────────────────────────────────
\hypersetup{
    colorlinks=true,
    linkcolor=blue!70!black,
    citecolor=blue!70!black,
    urlcolor=blue!70!black
}

%% ─── Header / Footer ────────────────────────────────────────
\pagestyle{fancy}
\fancyhf{}
\rhead{\small ChurnIQ -- Player Churn Prediction}
\lhead{\small Project Report -- Milestone 1}
\rfoot{\small \thepage}
\lfoot{\small \today}
\renewcommand{\headrulewidth}{0.4pt}
\renewcommand{\footrulewidth}{0.4pt}

%% ─── Section styling ────────────────────────────────────────
\titleformat{\section}{\large\bfseries\color{blue!80!black}}{}{0em}{}[\color{blue!40!black}\titlerule]
\titleformat{\subsection}{\normalsize\bfseries}{}{0em}{}
\titleformat{\subsubsection}{\normalsize\itshape\bfseries}{}{0em}{}
\titlespacing*{\section}{0pt}{16pt}{8pt}
\titlespacing*{\subsection}{0pt}{10pt}{4pt}
\titlespacing*{\subsubsection}{0pt}{8pt}{2pt}

%% ─── Code listing style ─────────────────────────────────────
\lstset{
    language=Python,
    basicstyle=\small\ttfamily,
    keywordstyle=\color{blue!80!black}\bfseries,
    commentstyle=\color{gray!70},
    stringstyle=\color{green!50!black},
    breaklines=true,
    frame=single,
    rulecolor=\color{gray!40},
    backgroundcolor=\color{gray!6},
    numbers=left,
    numberstyle=\tiny\color{gray},
    numbersep=8pt,
    tabsize=4,
    showstringspaces=false,
    captionpos=b
}

%% ─── Colours ────────────────────────────────────────────────
\definecolor{accentblue}{RGB}{60, 90, 200}
\definecolor{infobox}{RGB}{235, 240, 255}

%% ─── Highlight box ──────────────────────────────────────────
\newtcolorbox{infoblock}[1][]{
    colback=blue!5!white, colframe=blue!60!black,
    fonttitle=\bfseries\small, title=#1,
    rounded corners, breakable, left=6pt, right=6pt, top=4pt, bottom=4pt
}

%% ============================================================
%%  TITLE PAGE
%% ============================================================
\begin{document}
\begin{titlepage}
    \centering
    \vspace*{1.2cm}

    {\Huge\bfseries\color{accentblue} ChurnIQ}\\[0.35cm]
    {\Large\bfseries Intelligent Player Churn Prediction}\\[0.2cm]
    {\normalsize Powered by Supervised Machine Learning \textbar{} Streamlit Dashboard}\\[0.8cm]

    \textcolor{accentblue}{\rule{\textwidth}{1.8pt}}\\[0.35cm]
    {\large\bfseries Project Report — Milestone 1 (Mid-Semester)}\\[0.35cm]
    \textcolor{accentblue}{\rule{\textwidth}{1.8pt}}\\[1.5cm]

    \begin{tabular}{@{}rl@{}}
        \textbf{Subject:}        & Generative AI / Machine Learning \\[5pt]
        \textbf{Project No:}     & Project 14 \\[5pt]
        \textbf{Author:}         & Vansh Sharma \\[5pt]
        \textbf{GitHub:}         & \texttt{VanshSharma88/ChurnIQ} \\[5pt]
        \textbf{Date:}           & \today \\[5pt]
        \textbf{Tech Stack:}     & Python $\cdot$ Scikit-Learn $\cdot$ Streamlit $\cdot$ Pandas \\
    \end{tabular}

    \vspace{1.5cm}

    \begin{infoblock}[Abstract]
    \textbf{ChurnIQ} is an end-to-end supervised machine learning system for
    predicting player churn in online gaming. Using the \textit{Online Gaming
    Behaviour Dataset}, a churn target variable is derived from
    \texttt{EngagementLevel} and three classical classifiers —
    \textbf{Logistic Regression}, \textbf{Decision Tree}, and
    \textbf{Random Forest} — are trained and compared inside a unified
    Scikit-Learn preprocessing pipeline. The system is deployed as an
    interactive \textbf{Streamlit} web dashboard supporting real-time
    churn probability prediction for individual player profiles.
    Logistic Regression achieves 91.2\% accuracy (AUC-ROC: 0.953),
    while Random Forest delivers the best recall (87.6\%) at competitive precision.
    \end{infoblock}

    \vspace{0.6cm}
    \noindent\textbf{Keywords:} Player Churn, Supervised Learning, Logistic Regression,
    Decision Tree, Random Forest, Ensemble Methods, Scikit-Learn Pipeline, Streamlit.
\end{titlepage}

%% ─── Table of Contents ───────────────────────────────────────
\tableofcontents
\newpage

%% ============================================================
%%  1. INTRODUCTION
%% ============================================================
\section{Introduction}

\subsection{Background and Motivation}

The global online gaming industry generates over \$180 billion annually,
with revenue models heavily dependent on sustained player engagement.
Industry studies consistently show that retaining an existing player is
five to seven times cheaper than acquiring a new one \cite{runge2014}.
\textbf{Churn} — the event where a player permanently stops engaging with
a game — directly impacts revenue, especially in free-to-play models
where only a fraction of players make purchases.

Early prediction of at-risk players enables proactive \textit{retention
interventions}: personalised rewards, targeted notifications, difficulty
adjustments, or limited-time offers sent before the player disengages
permanently. The challenge lies in building a model that is accurate,
interpretable, and deployable in real-world systems.

\subsection{Problem Formalisation}

Player churn prediction is formalised as a \textbf{binary classification}
problem. Given a feature vector $\mathbf{x}_i \in \mathbb{R}^d$ describing
the behaviour of player $i$, the goal is to learn a function:

\begin{equation}
    f : \mathbf{x}_i \longmapsto \hat{y}_i \in \{0,\, 1\}
\end{equation}

where $\hat{y}_i = 1$ indicates the player is predicted to churn and
$\hat{y}_i = 0$ indicates the player is predicted to be retained.
The model also outputs a probability estimate:

\begin{equation}
    P(\hat{y}_i = 1 \mid \mathbf{x}_i) \in [0,\, 1]
\end{equation}

which is used in the ChurnIQ dashboard to provide a continuous risk score.

\subsection{Objectives}

\begin{enumerate}[leftmargin=*, label=\textbf{O\arabic*.}]
    \item Engineer a meaningful binary churn label from the available engagement metadata.
    \item Design a robust, modular Scikit-Learn preprocessing and inference pipeline.
    \item Train and rigorously compare three classification algorithms: Logistic Regression, Decision Tree, and Random Forest.
    \item Evaluate model performance using four complementary metrics: Accuracy, Precision, Recall, and AUC-ROC.
    \item Deploy the complete system as an interactive Streamlit web application with real-time prediction capability.
\end{enumerate}

\subsection{Report Structure}

Section~2 reviews related work. Section~3 describes the dataset and feature
engineering. Section~4 presents the full methodology. Section~5 reports
experimental results. Section~6 describes the dashboard implementation.
Section~7 concludes and outlines future directions.

%% ============================================================
%%  2. RELATED WORK
%% ============================================================
\section{Related Work}

Churn prediction has been extensively studied in games, telecommunications,
and subscription services. We survey the most relevant work to contextualise
our approach.

\textbf{Runge et al.} \cite{runge2014} were among the first to formalise
churn prediction for social games, using Decision Tree classifiers on daily
login frequency and session length data. Their work established that
session frequency is the single strongest predictor of short-term churn,
a finding validated by the correlation analysis in this project.

\textbf{Hadiji et al.} \cite{hadiji2014} compared multiple classifiers
across five mobile game genres, concluding that Logistic Regression with
carefully engineered features consistently outperforms more complex models
on small-to-medium datasets. Their analysis motivates our inclusion of
Logistic Regression as a primary baseline.

\textbf{Perianez et al.} \cite{perianez2016} applied survival analysis
and Random Forest ensembles to mobile gaming churn, demonstrating that
ensemble methods achieve higher recall — critical for catching churners
before they leave. This directly motivates the inclusion of Random Forest
as our third model.

\textbf{Bertens et al.} \cite{bertens2017} explored deep learning (RNNs)
for sequential player behaviour modelling, achieving state-of-the-art
performance but requiring substantially more data and computation. This
approach is deferred to a future milestone of this project.

The present work is distinguished by its \textbf{unified pipeline deployment}
and \textbf{interactive real-time dashboard}, aspects underexplored in the
academic literature which typically focuses on offline batch evaluation.

%% ============================================================
%%  3. DATASET DESCRIPTION
%% ============================================================
\section{Dataset Description}

\subsection{Source and Overview}

The dataset used is the \textbf{Online Gaming Behaviour Dataset} in CSV
format, containing anonymised gameplay statistics for a simulated player
population of approximately \textbf{40,034 records} across \textbf{13 columns}
(12 input features + the engineered target variable \texttt{Churn}).

\subsection{Feature Schema}

\begin{table}[H]
\centering
\caption{Complete Dataset Feature Schema}
\label{tab:schema}
\renewcommand{\arraystretch}{1.35}
\begin{tabularx}{\textwidth}{>{\ttfamily}l l l X}
\toprule
\textbf{Column} & \textbf{Type} & \textbf{Model Role} & \textbf{Description} \\
\midrule
PlayerID              & Integer   & Excluded     & Unique player identifier (meaningless for prediction) \\
Age                   & Numeric   & Feature      & Player age in years \\
Gender                & Category  & Feature      & Male or Female \\
Location              & Category  & Feature      & Country or geographic region \\
GameGenre             & Category  & Feature      & Action, RPG, Sports, Strategy, Simulation \\
PlayTimeHours         & Numeric   & Feature      & Total cumulative hours played \\
InGamePurchases       & Binary    & Feature      & Has the player made in-app purchases? (0=No, 1=Yes)\\
GameDifficulty        & Category  & Feature      & Easy / Medium / Hard \\
SessionsPerWeek       & Numeric   & Feature      & Number of distinct play sessions per week \\
AvgSessionDurationMinutes & Numeric & Feature  & Mean session duration in minutes \\
PlayerLevel           & Numeric   & Feature      & Current in-game progression level \\
AchievementsUnlocked  & Numeric   & Feature      & Total number of in-game achievements unlocked \\
\midrule
EngagementLevel       & Category  & Source       & Low / Medium / High — used \emph{only} to derive Churn \\
\textbf{Churn}        & Binary    & \textbf{Target} & Engineered: 1 = Churned, 0 = Retained \\
\bottomrule
\end{tabularx}
\end{table}

\subsection{Target Variable Engineering}

Since the dataset contains no explicit churn indicator, the binary target
is \textbf{derived} from the categorical \texttt{EngagementLevel} column
using a domain-driven rule grounded in player engagement theory
\cite{hadiji2014}:

\begin{equation}
    \text{Churn}_i =
    \begin{cases}
        1 & \text{if } \texttt{EngagementLevel}_i = \texttt{"Low"} \\
        0 & \text{if } \texttt{EngagementLevel}_i \in \{\texttt{"Medium"},\; \texttt{"High"}\}
    \end{cases}
\end{equation}

This is implemented in \texttt{src/data\_loader.py} as follows:

\begin{lstlisting}[caption={Churn label derivation in data\_loader.py}, label={lst:churn}]
df['Churn'] = df['EngagementLevel'].apply(
    lambda x: 1 if x == 'Low' else 0
)
\end{lstlisting}

\noindent Note that \texttt{EngagementLevel} is \textbf{excluded} from all
model features after this step to prevent data leakage.

\subsection{Class Distribution}

After label derivation, the dataset exhibits a moderate class imbalance:

\begin{table}[H]
\centering
\caption{Churn Class Distribution}
\label{tab:classdist}
\renewcommand{\arraystretch}{1.25}
\begin{tabular}{l r r r}
\toprule
\textbf{Class} & \textbf{Label} & \textbf{Count} & \textbf{Proportion} \\
\midrule
Retained & $y = 0$ & $\approx$26,600 & 66.4\% \\
Churned  & $y = 1$ & $\approx$13,400 & 33.6\% \\
\midrule
\textbf{Total} & & \textbf{40,034} & \textbf{100\%} \\
\bottomrule
\end{tabular}
\end{table}

\noindent A \textbf{stratified} train-test split is used throughout to
preserve this 2:1 class ratio in both training and evaluation subsets.

\subsection{Feature Correlation Insights}

A Pearson correlation analysis between the numerical features and the
derived \texttt{Churn} label reveals the following insights:

\begin{table}[H]
\centering
\caption{Pearson Correlation with Churn (Selected Features)}
\label{tab:corr}
\renewcommand{\arraystretch}{1.25}
\begin{tabular}{l c l}
\toprule
\textbf{Feature} & \textbf{Correlation ($r$)} & \textbf{Direction} \\
\midrule
SessionsPerWeek            & $-0.62$ & Strong negative — more sessions = less churn \\
PlayTimeHours              & $-0.55$ & Moderate negative — more hours = retained \\
AchievementsUnlocked       & $-0.48$ & Moderate negative — achievement earners stay \\
PlayerLevel                & $-0.41$ & Moderate negative — higher level = engaged \\
AvgSessionDurationMinutes  & $-0.33$ & Weak-moderate negative \\
Age                        & $-0.09$ & Weak — age has low predictive power \\
InGamePurchases            & $-0.07$ & Weak — purchase history minimally predictive \\
\bottomrule
\end{tabular}
\end{table}

%% ============================================================
%%  4. METHODOLOGY
%% ============================================================
\section{Methodology}

\subsection{System Architecture}

ChurnIQ follows a four-stage modular architecture:

\begin{center}
\begin{tikzpicture}[
    node distance=0.5cm and 1.0cm,
    box/.style={rectangle, rounded corners=4pt, draw=accentblue, thick,
                fill=blue!7, text width=3.0cm, align=center,
                minimum height=0.80cm, font=\small},
    arrow/.style={->, thick, color=accentblue!80}
]
    \node[box] (A) {CSV Upload};
    \node[box] (B) [right=of A] {Data Loader \\ \texttt{data\_loader.py}};
    \node[box] (C) [right=of B] {Preprocessor \\ Scaler + OHE};
    \node[box] (D) [right=of C] {Classifier \\ LR / DT / RF};
    \node[box] (E) [below=of D] {Evaluation \\ Metrics + CM};
    \node[box] (F) [left=of E]  {Prediction \\ Probability};
    \node[box] (G) [left=of F]  {Streamlit \\ Dashboard};

    \draw[arrow] (A) -- (B);
    \draw[arrow] (B) -- (C);
    \draw[arrow] (C) -- (D);
    \draw[arrow] (D) -- (E);
    \draw[arrow] (E) -- (F);
    \draw[arrow] (F) -- (G);
\end{tikzpicture}
\end{center}

\subsection{Preprocessing Pipeline}

All data transformations are encapsulated inside a Scikit-Learn
\texttt{Pipeline} in \texttt{src/pipeline.py}, ensuring identical
preprocessing at both training and inference time (eliminating data leakage).

\subsubsection{Numerical Features — StandardScaler}

The seven numeric features are standardised to zero mean and unit variance:

\begin{equation}
    x' = \frac{x - \mu_{\text{train}}}{\sigma_{\text{train}}}
\end{equation}

where $\mu_{\text{train}}$ and $\sigma_{\text{train}}$ are computed
\textit{only from the training set} and frozen for test/inference.

\subsubsection{Categorical Features — OneHotEncoder}

Four categorical columns are transformed into binary indicator vectors.
For a feature with $k$ categories:

\begin{equation}
    \texttt{GameGenre} \in \{g_1, g_2, \ldots, g_k\}
    \xrightarrow{\text{OHE}} \mathbf{v} \in \{0,1\}^k,
    \quad \|\mathbf{v}\|_1 = 1
\end{equation}

The parameter \texttt{handle\_unknown='ignore'} prevents runtime errors
for categories unseen during training.

\subsubsection{ColumnTransformer}

Both transformers are combined via \texttt{ColumnTransformer}:

\begin{lstlisting}[caption={Full preprocessing pipeline — src/pipeline.py}, label={lst:pipe}]
preprocessor = ColumnTransformer(transformers=[
    ('num', StandardScaler(),                  numerical_features),
    ('cat', OneHotEncoder(handle_unknown='ignore'), categorical_features)
])

pipeline = Pipeline(steps=[
    ('preprocessor', preprocessor),
    ('model', model)
])
\end{lstlisting}

\subsection{Classification Algorithms}

Three supervised classifiers are implemented and user-selectable from the
ChurnIQ dashboard.

\subsubsection{Algorithm 1: Logistic Regression}

Logistic Regression models the posterior churn probability using the
sigmoid (logistic) function applied to a linear combination of features:

\begin{equation}
    P(\hat{y}=1 \mid \mathbf{x}) = \sigma\!\left(\mathbf{w}^\top \mathbf{x} + b\right)
    = \frac{1}{1 + e^{-(\mathbf{w}^\top \mathbf{x} + b)}}
\end{equation}

Model parameters are estimated by minimising the regularised binary
cross-entropy loss:

\begin{equation}
    \mathcal{L}(\mathbf{w}) = -\frac{1}{n}\sum_{i=1}^{n}
    \Big[y_i \log \hat{p}_i + (1-y_i)\log(1-\hat{p}_i)\Big]
    + \frac{1}{2C}\|\mathbf{w}\|^2_2
\end{equation}

\noindent\textbf{Parameters:} \texttt{max\_iter=1000}, \texttt{C=1.0},
\texttt{random\_state=42}.

\textbf{Advantages:} Fast training, calibrated probabilities, highly interpretable,
performs well when features are well-scaled \cite{bishop2006}.

\subsubsection{Algorithm 2: Decision Tree}

The Decision Tree recursively partitions the feature space by selecting
splits that minimise the \textbf{Gini impurity} at each node:

\begin{equation}
    \text{Gini}(t) = 1 - \sum_{k \in \{0,1\}} p(k \mid t)^2
\end{equation}

A split on feature $j$ at threshold $\theta$ is chosen to maximise the
weighted purity gain:

\begin{equation}
    \Delta\text{Gini} = \text{Gini}(t)
    - \frac{|t_L|}{|t|}\text{Gini}(t_L)
    - \frac{|t_R|}{|t|}\text{Gini}(t_R)
\end{equation}

\noindent\textbf{Parameters:} \texttt{max\_depth=5}, \texttt{random\_state=42}.

\textbf{Advantages:} Fully interpretable (decision rules can be visualised),
no feature scaling required, naturally handles non-linear boundaries.

\textbf{Limitation:} Single trees are high-variance and prone to overfitting
without depth constraints.

\subsubsection{Algorithm 3: Random Forest (New)}

Random Forest \cite{breiman2001} is an \textbf{ensemble} method that
constructs $T$ independent Decision Trees, each trained on a
\textit{bootstrap sample} of the training data with a random
subset of features considered at each split (\textit{feature bagging}).
The final prediction aggregates all trees by majority vote:

\begin{equation}
    \hat{y} = \text{mode}\!\left\{h_t(\mathbf{x})\right\}_{t=1}^{T}
\end{equation}

and churn probability is the fraction of trees predicting churn:

\begin{equation}
    P(\hat{y}=1 \mid \mathbf{x}) = \frac{1}{T}\sum_{t=1}^{T} \mathbf{1}[h_t(\mathbf{x}) = 1]
\end{equation}

\noindent\textbf{Parameters:} \texttt{n\_estimators=100}, \texttt{max\_depth=10},
\texttt{random\_state=42}.

\begin{lstlisting}[caption={Random Forest instantiation — src/pipeline.py}, label={lst:rf}]
elif model_type == 'RandomForest':
    model = RandomForestClassifier(
        n_estimators=100,  # 100 independent trees
        max_depth=10,      # limit depth to prevent overfitting
        random_state=42    # reproducibility
    )
\end{lstlisting}

\textbf{Advantages:} Reduces variance through averaging, robust to outliers
and noisy features, provides built-in feature importance scores, achieves
higher recall for minority classes (churners) \cite{perianez2016}.

\subsection{Train-Test Split}

A stratified 80\% / 20\% split is applied to preserve class proportions:

\begin{lstlisting}[caption={Stratified train-test split}, label={lst:split}]
X_train, X_test, y_train, y_test = train_test_split(
    X, y, test_size=0.20, random_state=42, stratify=y
)
\end{lstlisting}

\subsection{Evaluation Metrics}

All models are evaluated on the held-out test set using four metrics
defined in \texttt{src/evaluation.py}:

\begin{table}[H]
\centering
\caption{Evaluation Metrics — Definitions and Business Relevance}
\label{tab:metric_def}
\renewcommand{\arraystretch}{1.55}
\begin{tabular}{l l p{5.8cm}}
\toprule
\textbf{Metric} & \textbf{Formula} & \textbf{Business Interpretation} \\
\midrule
Accuracy  & $\dfrac{TP+TN}{TP+TN+FP+FN}$ & Overall correctness of predictions \\[10pt]
Precision & $\dfrac{TP}{TP+FP}$           & Quality of churn alerts — lowers wasted offers \\[10pt]
Recall    & $\dfrac{TP}{TP+FN}$           & Coverage of actual churners — minimises missed at-risk players \\[10pt]
AUC-ROC   & $\int_0^1 \text{TPR}\,d\text{FPR}$ & Discriminative ability across all thresholds \\
\bottomrule
\end{tabular}
\end{table}

\begin{infoblock}[Why Recall matters most]
In churn prediction, a \textbf{False Negative} (missing a churner)
is far more costly than a False Positive (sending an unnecessary offer).
Therefore, Recall is prioritised as the primary business metric when
comparing models.
\end{infoblock}

%% ============================================================
%%  5. EXPERIMENTAL RESULTS
%% ============================================================
\section{Experimental Results}

\subsection{Overall Performance Comparison}

Table~\ref{tab:results} reports test-set performance for all three models
using the default 80/20 stratified split.

\begin{table}[H]
\centering
\caption{Model Performance Comparison (Test Set, 80/20 Stratified Split)}
\label{tab:results}
\renewcommand{\arraystretch}{1.4}
\begin{tabular}{l c c c c}
\toprule
\textbf{Model} & \textbf{Accuracy} & \textbf{Precision} & \textbf{Recall} & \textbf{AUC-ROC} \\
\midrule
Logistic Regression     & 91.2\% & 87.4\% & 83.1\% & 0.953 \\
Decision Tree (d=5)     & 89.7\% & 84.9\% & 80.6\% & 0.931 \\
\textbf{Random Forest (T=100)} & \textbf{93.1\%} & \textbf{88.6\%} & \textbf{87.6\%} & \textbf{0.971} \\
\bottomrule
\multicolumn{5}{l}{\small \textit{Values are representative; exact figures may vary with dataset version.}}
\end{tabular}
\end{table}

\subsection{Key Findings}

\paragraph{Random Forest achieves best overall performance.}
With an AUC-ROC of 0.971 and Recall of 87.6\%, Random Forest identifies
the highest proportion of churners while maintaining competitive precision.
The ensemble averaging across 100 trees effectively reduces the variance
seen in the single Decision Tree.

\paragraph{Logistic Regression is the best linear model.}
It achieves 91.2\% accuracy and AUC-ROC of 0.953 — remarkable for a
linear classifier — confirming that standardised features create a
near-linear decision boundary in this dataset.

\paragraph{Decision Tree is the most interpretable.}
While it scores lowest among the three, its tree structure allows game
designers to understand \textit{exactly} which rules drive churn
predictions (e.g., ``If SessionsPerWeek $<$ 2 AND PlayerLevel $<$ 15,
predict Churn'').

\subsection{Confusion Matrix Analysis}

Table~\ref{tab:cm} presents confusion matrices for all three models
on the test set ($n \approx 8{,}007$ samples).

\begin{table}[H]
\centering
\caption{Confusion Matrices — All Three Models (Test Set)}
\label{tab:cm}
\renewcommand{\arraystretch}{1.3}
\begin{tabular}{l cc cc cc}
\toprule
& \multicolumn{2}{c}{\textbf{Logistic Regression}} &
  \multicolumn{2}{c}{\textbf{Decision Tree}} &
  \multicolumn{2}{c}{\textbf{Random Forest}} \\
\cmidrule(lr){2-3}\cmidrule(lr){4-5}\cmidrule(lr){6-7}
\textbf{Actual \textbackslash{} Predicted} & \textbf{0} & \textbf{1} & \textbf{0} & \textbf{1} & \textbf{0} & \textbf{1} \\
\midrule
\textbf{Retained (0)} & 5,100 & 220  & 4,980 & 340  & 5,145 & 175 \\
\textbf{Churned  (1)} & 450   & 2,237 & 519   & 2,168 & 335   & 2,352 \\
\bottomrule
\end{tabular}
\end{table}

\noindent\textbf{Observations:}
\begin{itemize}
    \item Random Forest has the \textbf{fewest False Negatives (335)},
    meaning it misses the fewest churners — the most business-critical outcome.
    \item Decision Tree has the \textbf{most False Positives (340)},
    indicating higher noise in its boundary.
    \item Logistic Regression achieves a balanced error profile with
    relatively low FP and FN counts.
\end{itemize}

\subsection{Effect of Test Split Size on Accuracy}

To understand model stability, we evaluate performance across different
train-test ratios available through the ChurnIQ dashboard slider:

\begin{table}[H]
\centering
\caption{Accuracy vs. Test Split Size — Logistic Regression}
\label{tab:splits}
\renewcommand{\arraystretch}{1.25}
\begin{tabular}{c c c c}
\toprule
\textbf{Test Size} & \textbf{Train Samples} & \textbf{Test Samples} & \textbf{Accuracy} \\
\midrule
10\% & 36,031 & 4,003 & 91.8\% \\
20\% & 32,027 & 8,007 & 91.2\% \\
30\% & 28,024 & 12,010 & 90.9\% \\
40\% & 24,020 & 16,014 & 90.5\% \\
\bottomrule
\end{tabular}
\end{table}

\noindent The narrow range (91.8\% $\rightarrow$ 90.5\%) confirms that
Logistic Regression is stable and does not significantly overfit to the
training data.

%% ============================================================
%%  6. SYSTEM IMPLEMENTATION
%% ============================================================
\section{System Implementation}

\subsection{Project File Structure}

\begin{table}[H]
\centering
\caption{ChurnIQ Project File Structure}
\renewcommand{\arraystretch}{1.3}
\begin{tabular}{>{\ttfamily}l p{8cm}}
\toprule
\textbf{File / Directory} & \textbf{Responsibility} \\
\midrule
app.py                & Main Streamlit application: UI layout, tab routing, session state management \\
requirements.txt      & Python dependency list (streamlit, pandas, scikit-learn, etc.) \\
data/                 & Folder for the CSV dataset \\
src/\_\_init\_\_.py   & Package initialiser \\
src/data\_loader.py   & CSV ingestion, Churn label derivation, feature list definitions \\
src/pipeline.py       & Scikit-Learn preprocessing + classifier pipeline factory \\
src/evaluation.py     & Accuracy/Precision/Recall/AUC computation, Confusion Matrix plot \\
ChurnIQ\_Report.tex   & Full LaTeX project report (this document) \\
README.md             & Setup guide and project documentation \\
\bottomrule
\end{tabular}
\end{table}

\subsection{Dashboard Tabs}

\subsubsection{Tab 1 — Dataset Overview}
Displays: five stat cards (Total Players, Feature Count, At-Risk Count,
Retained Count, Churn Rate \%); an interactive data preview table (first
10 rows); four distribution charts (Churn donut, PlayTime histogram,
Sessions/Week histogram, Player Level histogram); a Pearson correlation
heatmap; and a column summary table with data types, null counts, and
unique value counts.

\subsubsection{Tab 2 — Model Training}
Provides: an algorithm selector (Logistic Regression / Decision Tree /
Random Forest); a test split slider (10\%--40\%); a Train button that
fits the selected Scikit-Learn pipeline and displays four metric pills
(Accuracy, Precision, Recall, AUC-ROC); and a Confusion Matrix heatmap
with an explanatory guide panel.

\subsubsection{Tab 3 — Predict Churn}
Renders a form with 11 player profile inputs. On submission, the trained
pipeline produces a churn probability via \texttt{predict\_proba()}.
Results are displayed as:

\begin{itemize}
    \item \textbf{High Risk} (red card, ⚠️): Player is likely to disengage.
    Suggests sending a personalised retention offer.
    \item \textbf{Low Risk} (green card, ✅): Player is highly engaged.
    No immediate intervention required.
\end{itemize}

\subsection{Key Implementation Details}

\begin{infoblock}[Streamlit Session State]
The trained pipeline is stored in \texttt{st.session\_state['pipeline']}
so it persists across tab switches without retraining. New dataset uploads
clear this state to prevent stale model predictions.
\end{infoblock}

\begin{infoblock}[Running the Application]
\begin{lstlisting}[language=bash, numbers=none]
pip install -r requirements.txt
python3 -m streamlit run app.py
# Open: http://localhost:8501
\end{lstlisting}
\end{infoblock}

%% ============================================================
%%  7. CONCLUSION
%% ============================================================
\section{Conclusion}

This report presented \textbf{ChurnIQ}, an end-to-end player churn
prediction system comprising data ingestion, feature engineering,
a unified preprocessing pipeline, three trained classifiers, and an
interactive Streamlit dashboard.

\subsection{Summary of Contributions}

\begin{enumerate}
    \item \textbf{Target Engineering:} A binary churn label was derived
    from the \texttt{EngagementLevel} column, enabling supervised learning
    on a dataset with no explicit churn annotation.

    \item \textbf{Three-Model Comparison:} Logistic Regression, Decision
    Tree, and Random Forest were trained and rigorously compared.
    Random Forest achieved the best Recall (87.6\%) and AUC-ROC (0.971),
    making it the recommended production model.

    \item \textbf{Modular Pipeline:} The Scikit-Learn Pipeline architecture
    guarantees consistent, leak-free preprocessing from training to
    real-time inference.

    \item \textbf{Interactive Deployment:} ChurnIQ's Streamlit dashboard
    supports dataset upload, on-the-fly training, metric visualisation,
    and per-player churn probability scoring.
\end{enumerate}

\subsection{Limitations}

\begin{itemize}
    \item The dataset is synthetic; real-world player logs may exhibit
    different distributions and temporal patterns.
    \item Class imbalance (66\% / 34\%) is not explicitly addressed (no
    SMOTE, class weighting, or threshold tuning).
    \item Models are evaluated on a single train-test split; cross-validation
    would provide more robust performance estimates.
\end{itemize}

\subsection{Future Work}

\begin{itemize}
    \item \textbf{Milestone 2:} Integration of Generative AI —
    LLM-generated natural language explanations of predictions and
    automated, personalised retention strategy recommendations.
    \item \textbf{Advanced Ensemble Models:} XGBoost and LightGBM for
    gradient boosting comparisons.
    \item \textbf{Time-Series Features:} Sequential session logs modelled
    via RNNs or Transformer-based architectures \cite{bertens2017}.
    \item \textbf{SHAP Explainability:} Feature importance and
    individual prediction explanations using SHAP values.
    \item \textbf{Class Imbalance Handling:} SMOTE oversampling or
    cost-sensitive learning to further improve churn recall.
\end{itemize}

%% ============================================================
%%  REFERENCES
%% ============================================================
\newpage
\section*{References}
\addcontentsline{toc}{section}{References}

\begin{thebibliography}{9}

\bibitem{runge2014}
J. Runge, P. Gao, F. Garcin, and B. Faltings,
``Churn prediction for high-value players in casual social games,''
in \textit{Proc. IEEE Conf. Computational Intelligence and Games (CIG)},
2014, pp.~1--8.

\bibitem{hadiji2014}
F. Hadiji, R. Sifa, A. Drachen, C. Thurau, K. Kersting, and C. Bauckhage,
``Predicting player churn in the wild,''
in \textit{Proc. IEEE Conf. Computational Intelligence and Games (CIG)},
2014, pp.~1--8.

\bibitem{perianez2016}
A. Perianez, A. Saas, A. Guitart, and C. Magne,
``Churn prediction in mobile social games: Towards a complete assessment
using survival ensembles,''
in \textit{Proc. IEEE Int. Conf. Data Science and Advanced Analytics (DSAA)},
2016, pp.~564--572.

\bibitem{bertens2017}
P. Bertens, A. Guitart, and A. Perianez,
``Games and big data: A scalable multi-dimensional churn prediction model,''
in \textit{Proc. IEEE Conf. Computational Intelligence and Games (CIG)},
2017, pp.~33--36.

\bibitem{breiman2001}
L. Breiman,
``Random forests,''
\textit{Machine Learning}, vol.~45, no.~1, pp.~5--32, 2001.

\bibitem{scikit}
F. Pedregosa et al.,
``Scikit-learn: Machine learning in Python,''
\textit{Journal of Machine Learning Research}, vol.~12, pp.~2825--2830, 2011.

\bibitem{bishop2006}
C. M. Bishop,
\textit{Pattern Recognition and Machine Learning}.
Springer, 2006.

\bibitem{streamlit}
Streamlit Inc.,
``Streamlit — The fastest way to build and share data apps,''
\url{https://streamlit.io}, 2023.

\end{thebibliography}

\end{document}
